\documentclass[12pt,a4paper]{article}

\usepackage[utf8]{inputenc}
\usepackage[T2A]{fontenc}
\usepackage[russian]{babel}
\usepackage{amsmath,amssymb}
\usepackage{enumitem}
\usepackage{array}
\usepackage{longtable}
\usepackage{geometry}
\usepackage{hyperref}

\geometry{margin=2.2cm}
\setlist{nosep}

\title{Билет №39 \\ \small Технология разработки и сопровождения программ. Жизненный цикл программы. Этапы разработки, степень и пути их автоматизации. (ИТМО) \\ Технология программирования. Жизненный цикл программы. Основные этапы. Инструментальные средства поддержки. (СпБПУ)}
\author{Сериков Владислав Александрович}
\date{13 мая 2026}

\begin{document}

\maketitle
\tableofcontents
\newpage

\section{Основные понятия}

\textbf{Программа} --- формализованное описание алгоритмов и данных, предназначенное для исполнения вычислительной системой.

\textbf{Программное обеспечение, ПО} --- совокупность программ, данных, документации, конфигурационных файлов, инструкций по установке, эксплуатации и сопровождению, необходимых для решения задач пользователя или организации.

\textbf{Программный продукт} --- ПО, поставляемое потребителю как самостоятельный результат разработки, обычно включающее исполняемые компоненты, исходный код, документацию, тесты, средства установки и сопровождения.

\textbf{Технология программирования} --- совокупность методов, процессов, моделей, стандартов, инструментов и организационных приёмов, применяемых для систематической разработки, проверки, внедрения, эксплуатации и сопровождения ПО.

\textbf{Инженерия программного обеспечения} --- дисциплина, изучающая систематический, измеримый и управляемый подход к созданию, эксплуатации и сопровождению программных систем.

\textbf{Жизненный цикл программного обеспечения} --- совокупность состояний, процессов и работ, через которые проходит программная система от возникновения потребности в ней до вывода из эксплуатации.

Жизненный цикл включает не только написание кода, но и:
\begin{itemize}
    \item анализ предметной области;
    \item формулирование требований;
    \item проектирование архитектуры и компонентов;
    \item программирование;
    \item тестирование и верификацию;
    \item документирование;
    \item внедрение;
    \item эксплуатацию;
    \item сопровождение;
    \item модернизацию или снятие системы с эксплуатации.
\end{itemize}

\section{Цели технологии разработки ПО}

Основные цели технологии разработки ПО:
\begin{enumerate}
    \item \textbf{Управляемость}: процесс разработки должен быть планируемым, контролируемым и измеримым.
    \item \textbf{Качество}: ПО должно удовлетворять функциональным и нефункциональным требованиям.
    \item \textbf{Надёжность}: система должна корректно работать в заданных условиях.
    \item \textbf{Сопровождаемость}: ПО должно быть удобно исправлять, изменять и расширять.
    \item \textbf{Повторяемость}: успешные практики разработки должны воспроизводиться в разных проектах.
    \item \textbf{Снижение рисков}: ошибки требований, архитектуры и реализации должны выявляться как можно раньше.
    \item \textbf{Экономическая эффективность}: стоимость разработки и сопровождения должна быть оправдана полезностью продукта.
\end{enumerate}

\section{Жизненный цикл программы}

\subsection{Общая структура жизненного цикла}

Типовой жизненный цикл программного продукта можно представить следующим образом:

\[
\text{Идея}
\rightarrow
\text{Требования}
\rightarrow
\text{Проектирование}
\rightarrow
\text{Реализация}
\rightarrow
\text{Тестирование}
\rightarrow
\text{Внедрение}
\rightarrow
\text{Эксплуатация}
\rightarrow
\text{Сопровождение}
\rightarrow
\text{Вывод из эксплуатации}.
\]

На практике эти стадии могут выполняться последовательно, итерационно, инкрементально или параллельно.

\subsection{Основные процессы жизненного цикла}

В жизненном цикле ПО обычно выделяют следующие группы процессов.

\subsubsection{Процессы соглашения}

К ним относятся:
\begin{itemize}
    \item приобретение ПО;
    \item поставка ПО;
    \item заключение договоров;
    \item согласование требований, сроков, бюджета и критериев приёмки.
\end{itemize}

\subsubsection{Организационные процессы}

Они обеспечивают возможность выполнения проекта:
\begin{itemize}
    \item управление проектом;
    \item управление качеством;
    \item управление рисками;
    \item управление конфигурацией;
    \item управление персоналом;
    \item совершенствование процессов разработки.
\end{itemize}

\subsubsection{Технические процессы}

К техническим процессам относятся:
\begin{itemize}
    \item анализ требований;
    \item проектирование архитектуры;
    \item детальное проектирование;
    \item программирование;
    \item интеграция;
    \item тестирование;
    \item верификация;
    \item валидация;
    \item установка;
    \item эксплуатация;
    \item сопровождение;
    \item утилизация или вывод из эксплуатации.
\end{itemize}

\section{Этапы разработки программного обеспечения}

\subsection{Постановка задачи}

На этом этапе формулируется исходная проблема, которую должна решать система.

Основные действия:
\begin{enumerate}
    \item определить заинтересованные стороны;
    \item описать проблему;
    \item определить цели системы;
    \item выделить ограничения;
    \item оценить осуществимость проекта.
\end{enumerate}

Результаты этапа:
\begin{itemize}
    \item описание предметной области;
    \item технико-экономическое обоснование;
    \item предварительное описание требований;
    \item решение о начале проекта.
\end{itemize}

\subsection{Анализ требований}

\textbf{Требование} --- условие или возможность, необходимая пользователю для решения задачи или системе для выполнения контракта, стандарта или спецификации.

Требования делятся на:
\begin{itemize}
    \item \textbf{функциональные} --- что система должна делать;
    \item \textbf{нефункциональные} --- какими свойствами должна обладать система;
    \item \textbf{ограничения} --- технологические, правовые, организационные или аппаратные рамки.
\end{itemize}

Примеры функциональных требований:
\begin{itemize}
    \item пользователь должен иметь возможность зарегистрироваться;
    \item система должна рассчитывать стоимость заказа;
    \item администратор должен иметь возможность блокировать учётные записи.
\end{itemize}

Примеры нефункциональных требований:
\begin{itemize}
    \item время ответа не более $200$ мс при $1000$ одновременных пользователей;
    \item доступность системы не ниже $99{,}9\%$;
    \item хранение паролей только в виде криптографических хешей.
\end{itemize}

Требования должны быть:
\begin{enumerate}
    \item полными;
    \item непротиворечивыми;
    \item проверяемыми;
    \item однозначными;
    \item трассируемыми;
    \item реализуемыми;
    \item актуальными.
\end{enumerate}

Результатом анализа требований является спецификация требований.

\subsection{Проектирование}

\textbf{Проектирование ПО} --- процесс преобразования требований в архитектуру, структуру данных, интерфейсы и алгоритмы будущей системы.

Проектирование делится на:
\begin{itemize}
    \item \textbf{архитектурное проектирование};
    \item \textbf{детальное проектирование}.
\end{itemize}

\subsubsection{Архитектурное проектирование}

Архитектура определяет крупные компоненты системы и связи между ними.

Архитектурное проектирование отвечает на вопросы:
\begin{itemize}
    \item из каких подсистем состоит система;
    \item какие интерфейсы имеют подсистемы;
    \item как распределяются данные;
    \item какие технологии используются;
    \item как обеспечиваются надёжность, безопасность, масштабируемость и производительность.
\end{itemize}

Примеры архитектурных стилей:
\begin{itemize}
    \item монолитная архитектура;
    \item клиент-серверная архитектура;
    \item многоуровневая архитектура;
    \item микросервисная архитектура;
    \item событийно-ориентированная архитектура;
    \item сервис-ориентированная архитектура.
\end{itemize}

\subsubsection{Детальное проектирование}

На этом уровне уточняются:
\begin{itemize}
    \item классы;
    \item функции;
    \item структуры данных;
    \item алгоритмы;
    \item форматы сообщений;
    \item схемы баз данных;
    \item обработка исключительных ситуаций.
\end{itemize}

Результаты проектирования:
\begin{itemize}
    \item архитектурная документация;
    \item диаграммы компонентов;
    \item диаграммы классов;
    \item диаграммы последовательностей;
    \item схемы данных;
    \item описание интерфейсов;
    \item прототипы.
\end{itemize}

\subsection{Реализация}

\textbf{Реализация} --- этап преобразования проектных решений в исходный код, конфигурации, скрипты сборки и исполняемые компоненты.

На этапе реализации выполняются:
\begin{itemize}
    \item выбор языка программирования;
    \item написание исходного кода;
    \item разработка модулей;
    \item написание модульных тестов;
    \item статический анализ кода;
    \item ревью кода;
    \item сборка приложения;
    \item подготовка документации разработчика.
\end{itemize}

Качественный код должен обладать:
\begin{itemize}
    \item понятностью;
    \item модульностью;
    \item низкой связностью;
    \item высокой связностью внутри модуля;
    \item тестируемостью;
    \item расширяемостью;
    \item соответствием стилю кодирования.
\end{itemize}

\subsection{Тестирование}

\textbf{Тестирование} --- процесс выполнения программы или анализа её компонентов с целью обнаружения дефектов и проверки соответствия требованиям.

Тестирование не доказывает отсутствие ошибок, но повышает уверенность в качестве системы.

Основные уровни тестирования:
\begin{enumerate}
    \item \textbf{Модульное тестирование} --- проверка отдельных функций, классов или модулей.
    \item \textbf{Интеграционное тестирование} --- проверка взаимодействия компонентов.
    \item \textbf{Системное тестирование} --- проверка всей системы.
    \item \textbf{Приёмочное тестирование} --- проверка системы заказчиком или пользователем.
\end{enumerate}

Основные виды тестирования:
\begin{itemize}
    \item функциональное;
    \item регрессионное;
    \item нагрузочное;
    \item стрессовое;
    \item тестирование безопасности;
    \item тестирование удобства использования;
    \item тестирование совместимости;
    \item тестирование восстановления после сбоев.
\end{itemize}

\subsection{Верификация и валидация}

\textbf{Верификация} отвечает на вопрос:
\[
\text{``Правильно ли мы строим продукт?''}
\]

Она проверяет соответствие промежуточных результатов спецификациям, стандартам и проектным решениям.

\textbf{Валидация} отвечает на вопрос:
\[
\text{``Правильный ли продукт мы строим?''}
\]

Она проверяет соответствие продукта реальным потребностям пользователей.

Примеры верификации:
\begin{itemize}
    \item инспекция требований;
    \item ревью архитектуры;
    \item статический анализ кода;
    \item проверка соответствия стандартам.
\end{itemize}

Примеры валидации:
\begin{itemize}
    \item демонстрация прототипа пользователю;
    \item приёмочные испытания;
    \item бета-тестирование;
    \item пользовательское тестирование.
\end{itemize}

\subsection{Внедрение}

\textbf{Внедрение} --- процесс передачи программной системы в эксплуатацию.

Внедрение включает:
\begin{itemize}
    \item установку ПО;
    \item настройку окружения;
    \item миграцию данных;
    \item обучение пользователей;
    \item настройку мониторинга;
    \item подготовку эксплуатационной документации;
    \item запуск в промышленную эксплуатацию.
\end{itemize}

Стратегии внедрения:
\begin{enumerate}
    \item \textbf{Прямое внедрение}: старая система сразу заменяется новой.
    \item \textbf{Параллельное внедрение}: старая и новая системы некоторое время работают одновременно.
    \item \textbf{Пилотное внедрение}: система сначала запускается на ограниченной группе пользователей.
    \item \textbf{Поэтапное внедрение}: функциональность вводится постепенно.
    \item \textbf{Blue-Green Deployment}: поддерживаются две среды, одна активная и одна резервная.
    \item \textbf{Canary Release}: новая версия сначала доступна малой доле пользователей.
\end{enumerate}

\subsection{Эксплуатация}

\textbf{Эксплуатация} --- стадия использования программной системы по назначению.

На этом этапе выполняются:
\begin{itemize}
    \item мониторинг работоспособности;
    \item резервное копирование;
    \item управление доступом;
    \item обработка инцидентов;
    \item анализ журналов;
    \item настройка производительности;
    \item обновление зависимостей;
    \item обеспечение безопасности.
\end{itemize}

\subsection{Сопровождение}

\textbf{Сопровождение ПО} --- процесс изменения программного продукта после поставки с целью исправления дефектов, улучшения характеристик, адаптации к изменившейся среде или предотвращения будущих проблем.

Выделяют четыре основных вида сопровождения:
\begin{enumerate}
    \item \textbf{Корректирующее сопровождение} --- исправление обнаруженных ошибок.
    \item \textbf{Адаптивное сопровождение} --- адаптация к изменениям внешней среды: ОС, СУБД, API, законодательства, оборудования.
    \item \textbf{Совершенствующее сопровождение} --- улучшение функциональности, производительности, удобства использования.
    \item \textbf{Профилактическое сопровождение} --- изменения, направленные на снижение будущих затрат и рисков: рефакторинг, обновление библиотек, улучшение тестов.
\end{enumerate}

Сопровождение часто занимает большую часть полной стоимости владения программной системой.

\subsection{Вывод из эксплуатации}

Вывод из эксплуатации включает:
\begin{itemize}
    \item прекращение поддержки;
    \item миграцию данных;
    \item архивирование;
    \item отключение сервисов;
    \item удаление устаревших компонентов;
    \item информирование пользователей;
    \item обеспечение юридических и регуляторных требований.
\end{itemize}

\section{Модели жизненного цикла}

\subsection{Каскадная модель}

\textbf{Каскадная модель} предполагает последовательное выполнение этапов:

\[
\text{Требования}
\rightarrow
\text{Проектирование}
\rightarrow
\text{Кодирование}
\rightarrow
\text{Тестирование}
\rightarrow
\text{Внедрение}
\rightarrow
\text{Сопровождение}.
\]

Преимущества:
\begin{itemize}
    \item простота планирования;
    \item понятная документация;
    \item удобство контроля этапов;
    \item применимость при стабильных требованиях.
\end{itemize}

Недостатки:
\begin{itemize}
    \item позднее обнаружение ошибок требований;
    \item сложность возврата к предыдущим этапам;
    \item слабая пригодность для проектов с изменяющимися требованиями;
    \item пользователь видит результат слишком поздно.
\end{itemize}

Каскадная модель подходит для проектов, где требования хорошо известны заранее, а цена изменений высока.

\subsection{V-образная модель}

V-образная модель уточняет связь между этапами разработки и соответствующими уровнями тестирования.

\[
\begin{array}{ccc}
\text{Требования} & & \text{Приёмочные тесты} \\
\text{Системное проектирование} & & \text{Системные тесты} \\
\text{Архитектура} & & \text{Интеграционные тесты} \\
\text{Детальный проект} & & \text{Модульные тесты} \\
& \text{Кодирование} &
\end{array}
\]

Идея модели: для каждого этапа проектирования заранее определяется соответствующий способ проверки.

\subsection{Итерационная модель}

В итерационной модели система создаётся через последовательность итераций. Каждая итерация включает анализ, проектирование, реализацию и тестирование.

\[
P_0 \rightarrow P_1 \rightarrow P_2 \rightarrow \dots \rightarrow P_n,
\]

где $P_i$ --- версия продукта после $i$-й итерации.

Преимущества:
\begin{itemize}
    \item раннее получение работающей версии;
    \item возможность уточнять требования;
    \item снижение рисков;
    \item регулярная обратная связь.
\end{itemize}

\subsection{Инкрементальная модель}

В инкрементальной модели продукт создаётся частями, называемыми инкрементами.

Каждый инкремент добавляет новую функциональность:

\[
S = I_1 \cup I_2 \cup \dots \cup I_n,
\]

где $S$ --- итоговая система, $I_i$ --- отдельный функциональный инкремент.

Минимальный пример:
\begin{itemize}
    \item инкремент 1: регистрация пользователей;
    \item инкремент 2: каталог товаров;
    \item инкремент 3: корзина;
    \item инкремент 4: оплата;
    \item инкремент 5: личный кабинет.
\end{itemize}

\subsection{Спиральная модель}

Спиральная модель ориентирована на управление рисками. Разработка проходит через витки спирали, каждый из которых включает:
\begin{enumerate}
    \item определение целей;
    \item анализ рисков;
    \item разработку и проверку решения;
    \item планирование следующего витка.
\end{enumerate}

Спиральная модель эффективна для крупных, дорогих и рискованных проектов.

\subsection{Agile-подходы}

\textbf{Agile} --- семейство гибких подходов, основанных на итерационной разработке, тесном взаимодействии с заказчиком и готовности к изменению требований.

Основные идеи:
\begin{itemize}
    \item работающий продукт важнее исчерпывающей документации;
    \item сотрудничество с заказчиком важнее формального согласования условий;
    \item готовность к изменениям важнее следования первоначальному плану;
    \item люди и взаимодействие важнее процессов и инструментов.
\end{itemize}

Распространённые Agile-подходы:
\begin{itemize}
    \item Scrum;
    \item Kanban;
    \item Extreme Programming;
    \item Lean Software Development.
\end{itemize}

\subsection{Scrum}

Scrum организует разработку короткими итерациями --- спринтами. Обычно спринт длится от одной до четырёх недель.

Основные роли:
\begin{itemize}
    \item \textbf{Product Owner} --- отвечает за ценность продукта и управляет backlog'ом;
    \item \textbf{Scrum Master} --- помогает команде соблюдать процесс и устранять препятствия;
    \item \textbf{Developers} --- создают продукт.
\end{itemize}

Основные артефакты:
\begin{itemize}
    \item Product Backlog;
    \item Sprint Backlog;
    \item Increment.
\end{itemize}

Основные события:
\begin{itemize}
    \item Sprint Planning;
    \item Daily Scrum;
    \item Sprint Review;
    \item Sprint Retrospective.
\end{itemize}

\subsection{DevOps}

\textbf{DevOps} --- совокупность практик, направленных на объединение разработки, эксплуатации и обеспечения качества для быстрой и надёжной поставки ПО.

DevOps включает:
\begin{itemize}
    \item автоматизацию сборки;
    \item автоматизацию тестирования;
    \item непрерывную интеграцию;
    \item непрерывную доставку;
    \item инфраструктуру как код;
    \item мониторинг;
    \item быструю обратную связь от эксплуатации к разработке.
\end{itemize}

\section{Степень и пути автоматизации этапов разработки}

\subsection{Общая идея автоматизации}

Автоматизация разработки ПО направлена на:
\begin{itemize}
    \item снижение числа ручных ошибок;
    \item ускорение повторяемых операций;
    \item повышение воспроизводимости результатов;
    \item улучшение контроля качества;
    \item обеспечение трассируемости изменений.
\end{itemize}

Не все работы можно полностью автоматизировать. Например, постановка целей, архитектурные компромиссы и переговоры с заказчиком требуют человеческого участия. Однако многие рутинные операции могут быть автоматизированы частично или полностью.

\subsection{Автоматизация анализа требований}

Автоматизируемые задачи:
\begin{itemize}
    \item ведение базы требований;
    \item трассировка требований к задачам, тестам и коммитам;
    \item контроль изменений требований;
    \item генерация отчётов;
    \item выявление противоречий и дубликатов с помощью средств анализа текста.
\end{itemize}

Инструменты:
\begin{itemize}
    \item системы управления требованиями;
    \item issue trackers;
    \item wiki-системы;
    \item системы управления документацией;
    \item средства построения UML/BPMN-моделей.
\end{itemize}

\subsection{Автоматизация проектирования}

Автоматизируемые задачи:
\begin{itemize}
    \item построение диаграмм;
    \item генерация каркаса кода по модели;
    \item обратное проектирование диаграмм по исходному коду;
    \item проверка архитектурных зависимостей;
    \item анализ связности и сложности компонентов.
\end{itemize}

Инструменты:
\begin{itemize}
    \item CASE-средства;
    \item UML-редакторы;
    \item средства моделирования баз данных;
    \item архитектурные анализаторы;
    \item генераторы кода.
\end{itemize}

\subsection{Автоматизация программирования}

Автоматизируемые задачи:
\begin{itemize}
    \item автодополнение кода;
    \item рефакторинг;
    \item форматирование;
    \item статический анализ;
    \item генерация шаблонного кода;
    \item управление зависимостями;
    \item сборка;
    \item контейнеризация.
\end{itemize}

Инструменты:
\begin{itemize}
    \item интегрированные среды разработки;
    \item компиляторы;
    \item интерпретаторы;
    \item системы сборки;
    \item пакетные менеджеры;
    \item линтеры;
    \item форматтеры;
    \item статические анализаторы;
    \item генераторы документации.
\end{itemize}

\subsection{Автоматизация тестирования}

Автоматизируемые задачи:
\begin{itemize}
    \item запуск модульных тестов;
    \item запуск интеграционных тестов;
    \item регрессионное тестирование;
    \item нагрузочное тестирование;
    \item тестирование API;
    \item проверка покрытия кода;
    \item генерация отчётов о тестировании.
\end{itemize}

Инструменты:
\begin{itemize}
    \item фреймворки модульного тестирования;
    \item средства mock/stub;
    \item инструменты нагрузочного тестирования;
    \item инструменты тестирования пользовательского интерфейса;
    \item системы контроля покрытия;
    \item CI/CD-платформы.
\end{itemize}

\subsection{Автоматизация внедрения}

Автоматизируемые задачи:
\begin{itemize}
    \item сборка релизного артефакта;
    \item упаковка в контейнер;
    \item развёртывание;
    \item настройка окружения;
    \item миграция схемы базы данных;
    \item откат версии;
    \item проверка работоспособности после релиза.
\end{itemize}

Инструменты:
\begin{itemize}
    \item CI/CD-системы;
    \item контейнерные платформы;
    \item системы оркестрации;
    \item инструменты Infrastructure as Code;
    \item системы управления конфигурациями.
\end{itemize}

\subsection{Автоматизация эксплуатации и сопровождения}

Автоматизируемые задачи:
\begin{itemize}
    \item мониторинг метрик;
    \item сбор логов;
    \item оповещение об инцидентах;
    \item автоматическое масштабирование;
    \item резервное копирование;
    \item обновление зависимостей;
    \item анализ уязвимостей;
    \item управление инцидентами.
\end{itemize}

Инструменты:
\begin{itemize}
    \item системы мониторинга;
    \item системы логирования;
    \item APM-системы;
    \item vulnerability scanners;
    \item системы управления инцидентами;
    \item системы резервного копирования.
\end{itemize}

\section{Инструментальные средства поддержки жизненного цикла}

\subsection{Системы управления версиями}

\textbf{Система управления версиями} хранит историю изменений файлов и позволяет нескольким разработчикам работать над проектом совместно.

Основные функции:
\begin{itemize}
    \item фиксация изменений;
    \item ветвление;
    \item слияние веток;
    \item откат к предыдущим версиям;
    \item просмотр истории;
    \item разрешение конфликтов;
    \item связывание изменений с задачами.
\end{itemize}

Пример минимального рабочего процесса:
\begin{enumerate}
    \item разработчик создаёт ветку для задачи;
    \item вносит изменения;
    \item запускает тесты;
    \item фиксирует изменения в коммите;
    \item отправляет ветку в общий репозиторий;
    \item создаёт запрос на слияние;
    \item проходит ревью;
    \item изменения попадают в основную ветку.
\end{enumerate}

\subsection{Системы управления задачами}

Они используются для:
\begin{itemize}
    \item планирования работ;
    \item назначения ответственных;
    \item отслеживания статусов;
    \item ведения backlog'а;
    \item контроля сроков;
    \item связи задач с требованиями, коммитами и тестами.
\end{itemize}

Типовые состояния задачи:
\[
\text{Open}
\rightarrow
\text{In Progress}
\rightarrow
\text{Code Review}
\rightarrow
\text{Testing}
\rightarrow
\text{Done}.
\]

\subsection{Средства непрерывной интеграции и доставки}

\textbf{Непрерывная интеграция, CI} --- практика частого объединения изменений в общую ветку с автоматическим запуском сборки и тестов.

\textbf{Непрерывная доставка, CD} --- практика, при которой система после прохождения автоматических проверок может быть быстро и надёжно поставлена в эксплуатационную среду.

\subsection{Средства статического анализа}

Статический анализ выполняется без запуска программы.

Он позволяет выявлять:
\begin{itemize}
    \item потенциальные ошибки;
    \item нарушения стиля кодирования;
    \item дублирование;
    \item недостижимый код;
    \item уязвимости;
    \item чрезмерную сложность;
    \item проблемы типизации.
\end{itemize}

\subsection{Средства динамического анализа}

Динамический анализ выполняется при запуске программы.

Он используется для:
\begin{itemize}
    \item поиска утечек памяти;
    \item профилирования производительности;
    \item анализа покрытия;
    \item выявления гонок данных;
    \item проверки поведения системы под нагрузкой.
\end{itemize}

\subsection{Средства документирования}

Документация бывает:
\begin{itemize}
    \item пользовательская;
    \item эксплуатационная;
    \item архитектурная;
    \item проектная;
    \item документация API;
    \item документация для разработчиков.
\end{itemize}

Средства документирования позволяют:
\begin{itemize}
    \item генерировать документацию из комментариев;
    \item вести базу знаний;
    \item поддерживать архитектурные решения;
    \item описывать API;
    \item публиковать документацию автоматически.
\end{itemize}

\section{Управление конфигурацией}

\textbf{Конфигурация ПО} --- совокупность версий программных компонентов, библиотек, настроек, документации, тестов и окружения, образующих конкретное состояние системы.

\textbf{Управление конфигурацией} --- деятельность по идентификации, контролю, учёту и проверке элементов конфигурации.

Основные задачи:
\begin{enumerate}
    \item определить элементы конфигурации;
    \item обеспечить версионирование;
    \item контролировать изменения;
    \item фиксировать состояние релизов;
    \item обеспечивать воспроизводимость сборки;
    \item поддерживать трассируемость.
\end{enumerate}

Элементы конфигурации:
\begin{itemize}
    \item исходный код;
    \item файлы сборки;
    \item конфигурационные файлы;
    \item документация;
    \item тестовые данные;
    \item схемы баз данных;
    \item контейнерные образы;
    \item скрипты миграции;
    \item внешние зависимости.
\end{itemize}

\section{Качество программного обеспечения}

\subsection{Понятие качества}

\textbf{Качество ПО} --- степень соответствия программного продукта установленным и предполагаемым потребностям пользователя.

Качество определяется не только отсутствием ошибок, но и удобством эксплуатации, сопровождения, развития и интеграции.

\subsection{Основные характеристики качества}

К важным характеристикам качества относятся:
\begin{enumerate}
    \item \textbf{Функциональная пригодность}: система реализует нужные функции.
    \item \textbf{Надёжность}: система сохраняет работоспособность при заданных условиях.
    \item \textbf{Производительность}: система эффективно использует ресурсы.
    \item \textbf{Удобство использования}: пользователю легко работать с системой.
    \item \textbf{Безопасность}: система защищает данные и операции.
    \item \textbf{Совместимость}: система взаимодействует с другими системами.
    \item \textbf{Сопровождаемость}: систему удобно изменять и анализировать.
    \item \textbf{Переносимость}: систему можно перенести в другую среду.
\end{enumerate}

\subsection{Метрики качества}

Примеры метрик:
\begin{itemize}
    \item количество дефектов на тысячу строк кода;
    \item покрытие кода тестами;
    \item цикломатическая сложность;
    \item среднее время восстановления после сбоя;
    \item среднее время между отказами;
    \item время ответа системы;
    \item количество уязвимостей;
    \item доля успешных сборок;
    \item длительность поставки изменения в production.
\end{itemize}

\section{Алгоритмы и типовые процедуры жизненного цикла}

\subsection{Алгоритм обработки изменения требования}

\textbf{Цель}: управляемо внести изменение в требования и связанные артефакты.

\textbf{Шаги алгоритма}:
\begin{enumerate}
    \item Зафиксировать запрос на изменение.
    \item Определить автора, причину и приоритет изменения.
    \item Найти связанные требования, модули, тесты и документы.
    \item Оценить влияние изменения на архитектуру, сроки, стоимость и риски.
    \item Принять решение: отклонить, отложить или принять изменение.
    \item Если изменение принято, обновить спецификацию требований.
    \item Создать задачи на реализацию.
    \item Реализовать изменение.
    \item Обновить тесты.
    \item Провести регрессионное тестирование.
    \item Обновить документацию.
    \item Зафиксировать изменение в журнале.
\end{enumerate}

\textbf{Схема}:
\[
\text{Запрос}
\rightarrow
\text{Анализ влияния}
\rightarrow
\text{Решение}
\rightarrow
\text{Реализация}
\rightarrow
\text{Проверка}
\rightarrow
\text{Документация}.
\]

\textbf{Минимальный пример}.

Было требование:
\[
R_1: \text{пользователь входит по логину и паролю}.
\]

Появилось новое требование:
\[
R_2: \text{добавить двухфакторную аутентификацию}.
\]

Анализ влияния показывает, что нужно изменить:
\begin{itemize}
    \item модуль авторизации;
    \item базу данных пользователей;
    \item пользовательский интерфейс входа;
    \item тесты безопасности;
    \item пользовательскую документацию.
\end{itemize}

\subsection{Алгоритм непрерывной интеграции}

\textbf{Цель}: автоматически проверять каждое изменение перед включением в основную ветку.

\textbf{Шаги алгоритма}:
\begin{enumerate}
    \item Разработчик вносит изменение в отдельной ветке.
    \item Изменение отправляется в удалённый репозиторий.
    \item CI-система получает событие о новом коммите.
    \item Загружается исходный код.
    \item Устанавливаются зависимости.
    \item Выполняется статический анализ.
    \item Выполняется сборка.
    \item Запускаются модульные тесты.
    \item Запускаются интеграционные тесты.
    \item Формируется отчёт.
    \item Если проверки успешны, изменение допускается к слиянию.
    \item Если проверки неуспешны, разработчик исправляет ошибки.
\end{enumerate}

\textbf{Схема}:
\[
\text{Commit}
\rightarrow
\text{Checkout}
\rightarrow
\text{Dependencies}
\rightarrow
\text{Static Analysis}
\rightarrow
\text{Build}
\rightarrow
\text{Tests}
\rightarrow
\text{Report}.
\]

\textbf{Минимальный пример}.

В проект добавлена функция:
\[
f(x) = x^2.
\]

Модульный тест:
\[
f(3) = 9.
\]

Если после изменения функция стала возвращать $x + x$, то тест для $x=3$ обнаружит ошибку:
\[
f(3) = 6 \neq 9.
\]

\subsection{Алгоритм выпуска релиза}

\textbf{Цель}: подготовить и безопасно поставить новую версию ПО.

\textbf{Шаги алгоритма}:
\begin{enumerate}
    \item Зафиксировать состав релиза.
    \item Назначить номер версии.
    \item Создать релизную ветку.
    \item Выполнить полную сборку.
    \item Выполнить регрессионное тестирование.
    \item Проверить миграции данных.
    \item Подготовить release notes.
    \item Создать релизный артефакт.
    \item Развернуть релиз в тестовой среде.
    \item Провести приёмочные проверки.
    \item Развернуть релиз в промышленной среде.
    \item Выполнить smoke-тесты.
    \item Настроить мониторинг.
    \item При критической ошибке выполнить откат.
\end{enumerate}

\textbf{Схема}:
\[
\text{Release branch}
\rightarrow
\text{Build}
\rightarrow
\text{Regression tests}
\rightarrow
\text{Deploy}
\rightarrow
\text{Smoke tests}
\rightarrow
\text{Monitoring}.
\]

\subsection{Алгоритм исправления дефекта}

\textbf{Цель}: воспроизвести, локализовать и устранить ошибку.

\textbf{Шаги алгоритма}:
\begin{enumerate}
    \item Зарегистрировать дефект.
    \item Описать ожидаемое и фактическое поведение.
    \item Указать окружение и шаги воспроизведения.
    \item Назначить приоритет и критичность.
    \item Воспроизвести ошибку.
    \item Локализовать причину.
    \item Написать тест, воспроизводящий дефект.
    \item Исправить код.
    \item Убедиться, что новый тест проходит.
    \item Запустить регрессионные тесты.
    \item Обновить документацию при необходимости.
    \item Закрыть дефект.
\end{enumerate}

\textbf{Минимальный пример}.

Ошибка: при делении $10$ на $0$ программа аварийно завершается.

Ожидаемое поведение: программа должна вернуть сообщение об ошибке.

Исправление:
\[
\text{если } b = 0,\ \text{то вернуть ошибку, иначе вернуть } a/b.
\]

\section{Документирование в жизненном цикле}

Документация является важным результатом почти каждого этапа жизненного цикла.

\subsection{Основные виды документации}

\begin{enumerate}
    \item \textbf{Документация требований}: описывает функции, ограничения и свойства системы.
    \item \textbf{Архитектурная документация}: описывает основные компоненты и их взаимодействие.
    \item \textbf{Проектная документация}: описывает внутреннюю структуру модулей.
    \item \textbf{Пользовательская документация}: объясняет, как пользоваться системой.
    \item \textbf{Эксплуатационная документация}: описывает установку, настройку, мониторинг и восстановление.
    \item \textbf{Документация API}: описывает интерфейсы программного взаимодействия.
    \item \textbf{Тестовая документация}: содержит планы тестирования, тест-кейсы, отчёты и данные.
\end{enumerate}

\subsection{Принципы хорошей документации}

Документация должна быть:
\begin{itemize}
    \item актуальной;
    \item полной;
    \item понятной;
    \item структурированной;
    \item проверяемой;
    \item связанной с артефактами проекта;
    \item удобной для поиска.
\end{itemize}

\section{Трассируемость}

\textbf{Трассируемость} --- возможность установить связи между требованиями, проектными решениями, кодом, тестами, дефектами и релизами.

Пример цепочки трассировки:
\[
\text{Требование}
\rightarrow
\text{Задача}
\rightarrow
\text{Коммит}
\rightarrow
\text{Тест}
\rightarrow
\text{Релиз}.
\]

Трассируемость помогает:
\begin{itemize}
    \item оценивать влияние изменений;
    \item доказывать выполнение требований;
    \item анализировать причины дефектов;
    \item проводить аудит;
    \item управлять качеством.
\end{itemize}

\section{Сравнение основных моделей жизненного цикла}

\begin{longtable}{|p{3.2cm}|p{4.2cm}|p{4.2cm}|p{3.2cm}|}
\hline
\textbf{Модель} & \textbf{Преимущества} & \textbf{Недостатки} & \textbf{Когда применять} \\
\hline
Каскадная &
Простота, строгая документация, понятные этапы &
Плохо переносит изменение требований &
Стабильные требования, регламентированные проекты \\
\hline
V-образная &
Чёткая связь разработки и тестирования &
Малая гибкость &
Критичные системы, где важна проверяемость \\
\hline
Итерационная &
Ранняя обратная связь, уточнение требований &
Сложнее планировать бюджет и сроки &
Проекты с не полностью ясными требованиями \\
\hline
Инкрементальная &
Ранние поставки полезной функциональности &
Требует хорошей архитектуры &
Продукты, которые можно поставлять частями \\
\hline
Спиральная &
Сильное управление рисками &
Высокая сложность управления &
Крупные и рискованные проекты \\
\hline
Agile &
Гибкость, быстрая обратная связь &
Требует зрелой команды и вовлечённого заказчика &
Динамичные продуктовые проекты \\
\hline
DevOps &
Быстрая и надёжная поставка, связь разработки и эксплуатации &
Требует высокой автоматизации и культуры взаимодействия &
Сервисы с частыми релизами \\
\hline
\end{longtable}

\section{Роль сопровождения в жизненном цикле}

Сопровождение является не второстепенным, а центральным процессом жизненного цикла. После первой поставки продукт обычно изменяется многократно.

Причины изменений:
\begin{itemize}
    \item обнаружение ошибок;
    \item изменение требований пользователей;
    \item изменение законодательства;
    \item изменение внешних API;
    \item появление новых угроз безопасности;
    \item рост нагрузки;
    \item устаревание технологий;
    \item необходимость интеграции с другими системами.
\end{itemize}

Для облегчения сопровождения применяются:
\begin{itemize}
    \item модульная архитектура;
    \item автоматические тесты;
    \item чистый код;
    \item документация;
    \item управление версиями;
    \item статический анализ;
    \item рефакторинг;
    \item мониторинг;
    \item логирование;
    \item управление конфигурацией.
\end{itemize}

\section{Рефакторинг}

\textbf{Рефакторинг} --- изменение внутренней структуры программы без изменения её внешнего поведения.

Цели рефакторинга:
\begin{itemize}
    \item улучшение читаемости;
    \item снижение сложности;
    \item устранение дублирования;
    \item повышение тестируемости;
    \item подготовка к добавлению новой функциональности.
\end{itemize}

Примеры рефакторингов:
\begin{itemize}
    \item переименование переменной;
    \item выделение функции;
    \item выделение класса;
    \item устранение дублирования;
    \item замена условной логики полиморфизмом;
    \item упрощение интерфейса.
\end{itemize}

Рефакторинг должен сопровождаться автоматическими тестами, иначе можно незаметно изменить поведение системы.

\section{Технический долг}

\textbf{Технический долг} --- накопленные в системе проектные, архитектурные и кодовые решения, которые ускорили разработку в краткосрочной перспективе, но увеличивают стоимость изменений в будущем.

Источники технического долга:
\begin{itemize}
    \item поспешные решения;
    \item недостаток тестов;
    \item устаревшие зависимости;
    \item слабая документация;
    \item дублирование кода;
    \item нарушение архитектуры;
    \item отсутствие автоматизации;
    \item недостаточный контроль качества.
\end{itemize}

Способы управления техническим долгом:
\begin{itemize}
    \item регулярный рефакторинг;
    \item архитектурные ревью;
    \item автоматические тесты;
    \item статический анализ;
    \item обновление зависимостей;
    \item выделение времени на улучшение качества;
    \item ведение списка технических задач.
\end{itemize}

\section{Связь этапов разработки и инструментов}

\begin{longtable}{|p{4cm}|p{6cm}|p{5cm}|}
\hline
\textbf{Этап} & \textbf{Основные работы} & \textbf{Инструментальная поддержка} \\
\hline
Постановка задачи &
Описание проблемы, целей и ограничений &
Документы, wiki, системы управления задачами \\
\hline
Анализ требований &
Сбор, анализ и согласование требований &
Requirements management, issue tracker, UML/BPMN \\
\hline
Проектирование &
Архитектура, интерфейсы, структура данных &
CASE-средства, UML, ER-моделирование \\
\hline
Реализация &
Написание кода, сборка, ревью &
IDE, компиляторы, VCS, build tools, linters \\
\hline
Тестирование &
Проверка корректности и качества &
Test frameworks, coverage tools, CI \\
\hline
Внедрение &
Развёртывание и настройка &
CI/CD, containers, IaC, orchestration \\
\hline
Эксплуатация &
Мониторинг, резервное копирование, инциденты &
Monitoring, logging, APM, backup tools \\
\hline
Сопровождение &
Исправления, адаптация, развитие &
VCS, issue tracker, CI/CD, static analysis \\
\hline
\end{longtable}

\section{Итоговые положения}

Технология разработки и сопровождения программ рассматривает создание ПО как управляемый инженерный процесс. Жизненный цикл программы охватывает путь от возникновения потребности до вывода системы из эксплуатации. Основные этапы включают анализ требований, проектирование, реализацию, тестирование, внедрение, эксплуатацию и сопровождение.

Современная разработка опирается на:
\begin{itemize}
    \item модели жизненного цикла;
    \item управление требованиями;
    \item архитектурное проектирование;
    \item автоматизированное тестирование;
    \item управление конфигурацией;
    \item непрерывную интеграцию и доставку;
    \item мониторинг и сопровождение;
    \item регулярное улучшение процессов.
\end{itemize}

Главная особенность программных систем состоит в том, что их разработка не заканчивается первым релизом. Большая часть работы связана с изменениями: исправлением ошибок, адаптацией к новым условиям, развитием функциональности и поддержанием качества. Поэтому технология программирования должна охватывать не только создание кода, но и весь жизненный цикл программного продукта.

\end{document}
